
\chapter{几何}





\subsection{2019年天台七下}

在以“探索光之奥秘”为主题的趣味物理实验中，用透明水箱模拟光线从空气射入某种液体，观察到入射角（\(\angle 1\)）与折射角（\(\angle 2\)）约为 \(4:3\) 的比例关系。为了挑战自我，同学们进一步思考：若两条入射光线以不同角度 \(\alpha\)、\(\beta\) 斜射入这种液体，液体内折射光线的夹角 \(\gamma\) 与 \(\alpha\)、\(\beta\) 的数学关系为（\hspace{1cm}）

\vspace{2mm}

\begin{center}
\includegraphics[width=0.4\textwidth]{2019天台七下01.png}
\end{center}

\noindent
\makebox[0.23\linewidth][l]{A. \(\alpha + \beta = \gamma\)}%
\hfill
\makebox[0.23\linewidth][l]{B. \(\dfrac{3}{4}(\alpha + \beta) = \gamma\)}%
\hfill
\makebox[0.23\linewidth][l]{C. \(\alpha + \beta = \dfrac{3}{4}\gamma\)}%
\hfill
\makebox[0.23\linewidth][l]{D. \(\dfrac{3}{4}(\alpha - \beta) = \gamma\)}

（答案见第~\pageref{几何-场景-题目1}~页）



\section{几何-比例}

\subsection{几何-比例-题目1}

\begin{center}
\includegraphics[width=1.0\textwidth]{几何-比例-题目1.png}
\end{center}

（答案见第~\pageref{几何-场景-题目1}~页）


\subsection{几何-比例-题目2}


\begin{center}
\includegraphics[width=1.0\textwidth]{几何-比例-题目2.png}
\end{center}

（答案见第~\pageref{几何-比例-题目2}~页）


\section{几何-求面积}
